\documentclass[12pt,a4paper]{article}
\usepackage[utf8]{inputenc}
\usepackage[brazilian]{babel}
\usepackage{amsmath}
\usepackage{amsfonts}
\usepackage{amssymb}
\usepackage{graphicx}
\usepackage{geometry}
\usepackage{hyperref}
\usepackage{listings}
\usepackage{xcolor}
\usepackage{booktabs}
\usepackage{float}

\geometry{margin=2.5cm}

\lstset{
    basicstyle=\ttfamily\small,
    breaklines=true,
    frame=single,
    language=Python,
    showstringspaces=false,
    commentstyle=\color{gray},
    keywordstyle=\color{blue},
    stringstyle=\color{red}
}

\title{\textbf{MAC0508 -- Introdução ao Processamento de Língua Natural}\\
\Large{EP1 - Parte B}}
\author{Leonardo Heidi Almeida Murakami \\ NUSP: 11260186}
\date{\today}

\begin{document}

\maketitle

\section{Descrição da Implementação}

\subsection{Arquitetura do Modelo}

Usamos um HMM (Hidden Markov Model) de ordem 2 para etiquetagem morfossintática (UPOS), seguindo a ideia de transições \(XY \to YZ\). O código está na classe \texttt{HMMTagger} (arquivo \texttt{hmm\_model.py}) e foi escrito de forma vetorizada com NumPy para melhorar a eficiência computacional, evitando loops Python desnecessários nas operações com matrizes.

O modelo \(\mu = (A, B, \pi)\) é definido assim:
\begin{itemize}
    \item \textbf{Estados}: as 17 etiquetas UPOS.
    \item \textbf{Transições} (matriz \(A\)): \(289\) linhas (pares de etiquetas \(XY\)) e \(17\) colunas (próxima etiqueta \(Z\)). Cada linha soma 1.
    \item \textbf{Emissões} (matriz \(B\)): \(17 \times |V|\), onde \(|V|\) é o vocabulário; linhas somam 1.
    \item \textbf{Iniciais} (vetor \(\pi\)): distribuição sobre as 17 etiquetas iniciais; soma 1.
\end{itemize}

\subsection{Pré-processamento}

O método \texttt{parse\_conllu()} lê os arquivos CoNLL-U e extrai:
\begin{itemize}
    \item Coluna 2 (FORM) como texto;
    \item Coluna 4 (UPOS) como etiqueta.
\end{itemize}
Linhas de comentário e linhas vazias são ignoradas, assim como tokens multi-palavra (IDs com `-'). Não adicionamos tokens especiais de início/fim de sentença nas palavras; trabalhamos diretamente com a sequência de etiquetas para o HMM.

\subsection{Inicialização dos Parâmetros}

As matrizes são inicializadas por contagem no córpus de treinamento, com \textbf{Lidstone} (\(\alpha = 1\)) para evitar zeros:
\begin{align*}
\pi_i &= \frac{\text{contagem}(S_i) + \alpha}{\sum_j \text{contagem}(S_j) + \alpha\,|S|} \\
A_{(X,Y),Z} &= \frac{\text{contagem}(XYZ) + \alpha}{\text{contagem}(XY) + \alpha\,|S|} \\
B_{i,w} &= \frac{\text{contagem}(S_i, w) + \alpha}{\text{contagem}(S_i) + \alpha\,|V|}
\end{align*}
onde \(|S|=17\).

\subsection{Palavras Desconhecidas}

Usamos o conjunto de validação para identificar palavras fora do vocabulário e adicionar uma coluna \texttt{<DESC>} em \(B\). Estimamos \(B_{i,\texttt{<DESC>}}\) a partir das etiquetas observadas quando a palavra é desconhecida no \textit{dev}, e renormalizamos as linhas de \(B\). Em execução, palavras não vistas mapeiam para \texttt{<DESC>}.

\subsection{Treinamento: Baum–Welch}

O treinamento opcional refina \(A\), \(B\) e \(\pi\) com Baum–Welch (\textit{EM}), usando \(\alpha\) de Lidstone em cada reestimação. O critério de parada compara a distância entre iterações:
\[
\mathrm{dist}(\mu^s, \mu^{s+1}) = \sum_{i,j}\big(A_{ij}^{s+1}-A_{ij}^s\big)^2 + \sum_{i,k}\big(B_{ik}^{s+1}-B_{ik}^s\big)^2 + \sum_i\big(\pi_i^{s+1}-\pi_i^s\big)^2,
\]
encerrando quando \(\mathrm{dist}(\mu^s,\mu^{s+1})/\mathrm{dist}(\mu^0,\mu^1) < \varepsilon\) (padrão \(\varepsilon=10^{-2}\)), ou após 100 iterações.

\subsection{Etiquetagem: Viterbi de ordem 2}

Para cada sentença, aplicamos Viterbi de ordem 2. A implementação usa um tensor de \(A\) (\(17\times 17\times 17\)) para acelerar as transições \(XY\to YZ\).

\section{Ativação do Programa}

\subsection{Requisitos}

\begin{itemize}
    \item Python 3.7 ou superior
    \item Bibliotecas: \texttt{numpy}, \texttt{tqdm}
\end{itemize}

Instalação rápida:
\begin{lstlisting}[language=bash]
pip install numpy tqdm
\end{lstlisting}

\subsection{Arquivos necessários}

Coloque estes arquivos no mesmo diretório:
\begin{itemize}
    \item \texttt{hmm\_model.py}
    \item \texttt{pt\_porttinari-ud-train.conllu}
    \item \texttt{pt\_porttinari-ud-dev.conllu}
    \item \texttt{pt\_porttinari-ud-test.conllu}
\end{itemize}

\subsection{Execução}

Para treinar e avaliar:
\begin{lstlisting}[language=bash]
python hmm_model.py
\end{lstlisting}

O script verifica os arquivos, treina o HMM (com Baum–Welch), etiqueta o teste com Viterbi e imprime as métricas.

\section{Resultados}

\subsection{Resumo do Treinamento}

\begin{table}[H]
\centering
\begin{tabular}{lr}
\toprule
\textbf{Métrica} & \textbf{Valor} \\
\midrule
Sentenças de treino & 5.893 \\
Número de etiquetas & 16 \\
Tamanho do vocabulário & 17.022 \\
Pares de estados (XY) & 256 \\
Palavras desconhecidas (dev) & 1.302 \\
Ocorrências desconhecidas (dev) & 1.362 \\
Iterações Baum--Welch & 2 \\
Razão dist. final/inicial & $2{,}69\times10^{-4}$ \\
\bottomrule
\end{tabular}
\caption{Resumo do treinamento (logs do programa).}
\end{table}

\subsection{Métricas Globais}

\begin{table}[H]
\centering
\begin{tabular}{lr}
\toprule
\textbf{Métrica} & \textbf{Valor} \\
\midrule
Acurácia & 56,72\% \\
Precisão Média & 72,61\% \\
Cobertura Média & 52,02\% \\
Medida-F (média) & 54,03\% \\
Medida-F (global) & 60,62\% \\
Etiquetas corretas & 19.053 / 33.594 \\
\bottomrule
\end{tabular}
\caption{Métricas globais na avaliação do conjunto de teste.}
\end{table}

\subsection{Top 10 etiquetas mais fáceis de prever}

\begin{table}[H]
\centering
\begin{tabular}{clrrrr}
\toprule
\textbf{\#} & \textbf{Etiqueta} & \textbf{F1} & \textbf{Precisão} & \textbf{Cobertura} & \textbf{Corretas} \\
\midrule
1 & SYM   & 92,62\% & 91,13\% & 94,17\% & 113 \\
2 & ADP   & 91,44\% & 99,14\% & 84,86\% & 4.259 \\
3 & AUX   & 88,74\% & 97,52\% & 81,40\% & 788 \\
4 & CCONJ & 88,56\% & 100,00\% & 79,47\% & 654 \\
5 & DET   & 82,65\% & 72,77\% & 95,63\% & 4.685 \\
6 & ADV   & 78,40\% & 88,77\% & 70,20\% & 893 \\
7 & PUNCT & 72,03\% & 100,00\% & 56,29\% & 2.501 \\
8 & PRON  & 68,49\% & 67,96\% & 69,02\% & 891 \\
9 & NUM   & 54,39\% & 100,00\% & 37,36\% & 229 \\
10 & VERB & 40,93\% & 93,41\% & 26,20\% & 893 \\
\bottomrule
\end{tabular}
\caption{Etiquetas com melhor desempenho (por F1).}
\end{table}

\subsection{Top 10 etiquetas mais difíceis de prever}

\begin{table}[H]
\centering
\begin{tabular}{clrrrr}
\toprule
\textbf{\#} & \textbf{Etiqueta} & \textbf{F1} & \textbf{Precisão} & \textbf{Cobertura} & \textbf{Corretas} \\
\midrule
1 & SCONJ & 19,92\% & 78,79\% & 11,40\% & 52 \\
2 & ADJ   & 21,98\% & 12,61\% & 85,50\% & 1.433 \\
3 & NOUN  & 30,78\% & 93,21\% & 18,43\% & 1.167 \\
4 & PROPN & 33,51\% & 66,44\% & 22,41\% & 495 \\
5 & VERB  & 40,93\% & 93,41\% & 26,20\% & 893 \\
6 & NUM   & 54,39\% & 100,00\% & 37,36\% & 229 \\
7 & PRON  & 68,49\% & 67,96\% & 69,02\% & 891 \\
8 & PUNCT & 72,03\% & 100,00\% & 56,29\% & 2.501 \\
9 & ADV   & 78,40\% & 88,77\% & 70,20\% & 893 \\
10 & DET  & 82,65\% & 72,77\% & 95,63\% & 4.685 \\
\bottomrule
\end{tabular}
\caption{Etiquetas com pior desempenho relativo (por F1).}
\end{table}

\subsection{Top 10 palavras mais difíceis de etiquetar}

\begin{table}[H]
\centering
\begin{tabular}{clrrr}
\toprule
\textbf{\#} & \textbf{Palavra} & \textbf{Taxa de erro} & \textbf{Acertos} & \textbf{Total} \\
\midrule
1 & .        & 100,00\% & 0 & 1.636 \\
2 & )        & 100,00\% & 0 & 128 \\
3 & mas      & 100,00\% & 0 & 61 \\
4 & :        & 100,00\% & 0 & 49 \\
5 & ?        & 100,00\% & 0 & 46 \\
6 & -        & 100,00\% & 0 & 44 \\
7 & Mas      & 100,00\% & 0 & 36 \\
8 & empresa  & 100,00\% & 0 & 29 \\
9 & Segundo  & 100,00\% & 0 & 27 \\
10 & Para    & 100,00\% & 0 & 26 \\
\bottomrule
\end{tabular}
\caption{Palavras com maior taxa de erro (apenas termos com frequência mínima de 3).}
\end{table}

\section{Comentários e Conclusões}

A acurácia de 56,72\% ficou abaixo do esperado. Embora o algoritmo de Baum--Welch tenha convergido rapidamente (apenas 2 iterações), os resultados mostram que o modelo tem dificuldades importantes.

Um ponto que chamou atenção foi a discrepância enorme no desempenho entre diferentes etiquetas. Enquanto ADP e SYM tiveram F1 acima de 90\%, outras como SCONJ (19,92\%) e ADJ (21,98\%) ficaram completamente perdidas. Isso sugere que o contexto de pares de etiquetas não é suficiente para disambiguar casos mais complexos.

O resultado mais estranho aparece quando olhamos as palavras mais difíceis de etiquetar. Apesar de PUNCT ter F1 de 72\% e precisão de 100\%, praticamente todas as pontuações específicas tiveram taxa de erro de 100\%: o ponto final errou todas as 1.636 ocorrências, o parêntese fechado errou as 128 vezes, e assim por diante. Isso não faz sentido à primeira vista.

Analisando com calma, percebe-se que a cobertura de PUNCT é apenas 56\%, ou seja, o modelo está deixando de identificar quase metade das pontuações. Quando ele prevê PUNCT, acerta (daí os 100\% de precisão), mas ele simplesmente não está prevendo PUNCT para os símbolos mais comuns como ponto, vírgula e parênteses. Provavelmente o modelo está confundindo essas marcações com outras etiquetas, talvez por viés das sequências de treinamento.

Outro problema visível é com palavras ambíguas como ``mas'' e ``Mas'', que também erraram 100\%. Essas palavras podem ser tanto conjunções (CCONJ) quanto advérbios (ADV) dependendo do contexto, e o HMM claramente não está capturando essas nuances.

\end{document}


